\section{日期总结}
\begin{enumerate}
  \item 日。基本都是5、7、10、15。
        \begin{enumerate}
          \item 5。执行保全届满、同时申请保全和支付令。
          \item 7。起诉受理、小额举证。
          \item 10。期间送达判决、执行异议、和其他部门交涉期、裁定上诉期。
          \item 15。举证期、主动方申请期、法院审查期、被动方异议期/宽限期。
        \end{enumerate}
  \item 月。当事人提申请基本是 6，法院审基本期是 3，仲裁救济审是 2。
  \item 年。基本是 1，执行 2，诉讼时效 3。
  \item 离婚。
        \begin{enumerate}
          \item If(一方离开超 1 年) + 原告住所地。
          \item 没离成 → 6 个月后、新情况、被告随时。
          \item 无效婚宣告申请。夫妻生前均可申请，夫妻一方 || 双方死亡 1 年内亦可。
        \end{enumerate}
  \item 举证。
        \begin{enumerate}
          \item 协商 || 法院确定的举证期限 \dydy 15 日。
          \item 简易程序举证期限、答辩期限协商 || 法院确定。≤ 15 日。
          \item 小额程序举证期限。协商 || 法院确定。≤ 7 日。
        \end{enumerate}
  \item 保全。
        \begin{enumerate}
          \item 必须 48 小时内裁定，30 日内起诉（保全手续移送受理法院）|| 仲裁，否则法院裁定解除保全。
          \item 执行保全。履行期限届满后 5 日内申请执行，否则法院解除保全。
          \item 妨碍执行。15 日以下拘留。
        \end{enumerate}
  \item 送达。
  \item \begin{enumerate}
          \item 期间的耽误。障碍消除后的 10 日内申请顺延，是否允许法院决定。
          \item 公告送达。兜底送达。国内 30 日，涉外 3 个月。
        \end{enumerate}
  \item 一审。
        \begin{enumerate}
          \item 提交答辩状期间/答辩期。一（异）坨大便。答辩、提管辖权异议。
          \item 举证期限届满前。跟证据和证明有关。
          \item 法庭辩论终结前。请求、回避。诉讼请求要定稿、申请回避。
          \item 起诉 7 日内处理。
          \item 审期 6 + 6。涉外无要求。
          \item 当庭宣判 10 日给，定期宣判立即给。
        \end{enumerate}
  \item 公益诉讼。
        \begin{enumerate}
          \item 初步证据即可 + 侵权行为地 || 被告住所地中院。时效 3 年。
          \item 法院应当在受理后 10 日内书面告知有关行政主管部门。
          \item 行政公益诉讼。起诉期 6 个月。
        \end{enumerate}
  \item 简易程序。审期 3 + 1。
  \item 小额。审期 2 + 1。
  \item 三撤。知道 || 应当知道 6 个月内向作出该生效/终审判决的法院申请。
  \item 二审。
        \begin{enumerate}
          \item 上诉期 → 判决 15 日，裁定 10 日。
          \item 审理 30 天或 3 个月。
        \end{enumerate}
  \item 再审。
        \begin{enumerate}
          \item 申请。自知道 || 应当知道裁判文书生效之日起的 6 个月内提出。
          \item 接受抗诉的法院应当在收到抗诉书之日起 30 日内直接裁定再审。
          \item 审查。一方找上级，一方找原审 → 无法协商 → 原审受理 → 3 月审查。
          \item 同一审。
        \end{enumerate}
  \item 特别程序。
        \begin{enumerate}
          \item 立案之日 || 公告之日起，30 日内审结；选民资格 → 选举开始前。
          \item 确认调解协议效力。生效之日起 30 日内共同提出申请。
          \item 异议期。当事人 15 日。②. 利害人知道 || 应当知道 6 个月。
        \end{enumerate}
  \item 支付令。
        \begin{enumerate}
          \item 已经申请诉前保全的不允许申请支付令 → 5 日内通知不予受理。
          \item 符合条件 → 15 日内发出支付令。作出之日起发生法律效力。
          \item 15 日内，债务人未清偿债务 || 未书面提出实质异议 → 强制执行。
        \end{enumerate}
  \item 公示催告。
        \begin{enumerate}
          \item 3 日内公告。公示催告期间 ≥ 60 日（法院决定）。期间转让票据行为无效。
          \item 自知道 || 应当知道判决公告之日起的 1 年内向作出除权判决的法院起诉。
        \end{enumerate}
  \item 执行。
        \begin{enumerate}
          \item 申请执行的时效 2 年，可以中止、中断。必须以生效法律文书为根据。
          \item 执行异议。收到执行通知之日起的 10 日内提。
          \item 10 日内、可以委托当地法院代为执行。
          \item 对执行标的的异议。15 日内审查。
          \item 代位申请执行债权。15 日内，次债务人未清偿债务 || 未书面 || 口头提出实质异议。
        \end{enumerate}
  \item 仲裁。
        \begin{enumerate}
          \item 仲裁裁决。边裁 1 票，首裁 1.5 票。作出之日生效。
          \item 撤销仲裁裁决。当事人收到裁决之日起 6 个月内向仲裁委所在地中院申请撤销/不予执行。
          \item 法院审查。2 月、合议庭。
        \end{enumerate}
\end{enumerate}

\section{找谁总结}
\begin{enumerate}
  \item 当事人、被动方、财产/标的物、机构/组织。
  \item 保全。向被申请人住所地、被保全财产/证据所在地、将来对案件有管辖权的法院申请。
  \item 执行保全。一审法院 || 与之同级的被执行财产所在地法院。
  \item 公益诉讼。侵权行为地 || 被告住所地中院。
  \item 三撤。作出该生效/终审判决的法院申请。
  \item 再审。原则找上级法院。一方人数众多 || 双方都是公民可以找原审法院。
  \item 确认调解。当事人住所地、标的物所在地、调解组织所在地法院（基层 \hz 中级）。
  \item 实现担保物权。担保财产所在地 || 担保物权登记地的基层法院管辖。海事的归海事法院。
  \item 支付令。向债务人住所地的基层法院提出申请。
  \item 公示催告程序。由票据最后持有人向票据付款地法院申请。
  \item 执行。一审法院 + 被执行财产所在地法院。
  \item 仲裁。
        \begin{enumerate}
          \item 确认效力。仲裁协议签订地、约定的仲裁机构所在地、申请人/被申请人住所地中院。
          \item 仲裁执行。被执行人住所地、被执行财产所在地的中院。
          \item 撤销、不予执行。仲裁委所在地中院。
        \end{enumerate}
\end{enumerate}
